%%%%%%%%%%  scan idx 69  |  folio 56  %%%%%%%%%%
\begin{equation}\tag{2.30}
 \underline K(Y)\longrightarrow\underline K(X)\longrightarrow\underline K(U)\longrightarrow 0
\end{equation}
\emph{est exacte.}

N.B. On comparera avec la proposition 2.3.

Le fait que $\underline K(X)\longrightarrow\underline K(U)$ soit surjectif résulte aussitôt de la proposition 2.5 (ou aussi de la proposition 2.6). Reste à prouver que si $x\in\underline K(X)$ a une restriction nulle à $\underline K(U)$, il est dans l'image de $\underline K(Y)$. Pour ceci, on revient à la définition explicite de $\underline K(X)$ et de $\underline K(U)$~: on a deux faisceaux algébriques cohérents $F$ et $G$ sur $X$ avec $x=\gamma_X(F)-\gamma_X(G)$ et $\gamma_U(F)-\gamma_U(G)=0$~; cette dernière relation signifie qu'on peut trouver des faisceaux cohérents $P'$ et $Q'$ sur $U$, munis de filtrations dont les facteurs soient deux à deux isomorphes sur $U$, et tels que l'on ait un isomorphisme sur $U$~:
\[
 P=F+Q'\longrightarrow Q=G+P'.
\]
En vertu de la proposition 2.5, les faisceaux $P'$ et $Q'$ sont des restrictions à $U$ de faisceaux cohérents sur $X$ (notés par les mêmes lettres)~; de même, leurs filtrations proviennent de filtrations sur les faisceaux correspondants sur $X$, en vertu du lemme 2.11 ci-dessous. On a alors~:
\[
 \gamma_X(F)-\gamma_X(G)=(\gamma_X(P)-\gamma_X(Q))+(\gamma_X(P')-\gamma_X(Q')).
\]
Introduisant les facteurs $P'_i$ et $Q'_j$ des filtrations de $P'$ et $Q'$ sur $X$, le dernier terme est la somme de termes $\gamma_X(P'_i)-\gamma_X(Q'_i)$, où $P'_i|_U\simeq Q'_i|_U$. On est ainsi ramené à prouver le cas particulier suivant~: si $R$ et $S$ sont des faisceaux cohérents sur $X$ tels que $R|_U\simeq S|_U$, alors $\gamma_X(R)-\gamma_X(S)$ est dans l'image de $\underline K(Y)$. Or soit $T$ le graphe de l'isomorphisme donné de $R|_U$ sur $S|_U$~; c'est un sous-faisceau cohérent de $(R\times S)|_U$, donc en vertu du lemme 2.11, c'est la restriction à $U$ d'un sous-faisceau cohérent, 
